\chapter{Logika modalna}\label{ch:ncl-modal}

\section{Język modalny i jego interpretacja}

W klasycznym rachunku zdań pytamy, czy zdanie jest
prawdziwe. W logice modalnej możemy dodatkowo pytać,
czy jest koniecznie prawdziwe lub czy jest możliwe.
Znaczenie tych pojęć zależy od zastosowania: mogą
dotyczyć możliwości logicznej, informacji posiadanych
przez podmiot albo dostępnych stanów pewnego układu.
W tym rozdziale badamy normalne zdaniowe logiki
modalne od strony syntaktycznej. Pełną semantykę
relacyjną oraz twierdzenia o pełności rozwiniemy
w następnym rozdziale.

\begin{definicja}[Język zdaniowej logiki modalnej]
\label{def:ncl-modal-language}
Niech $\mathrm{Var}$ będzie przeliczalnie nieskończonym
zbiorem zmiennych zdaniowych. Zbiór formuł modalnych
$\mathrm{Fm}_{\Box}$ definiujemy rekurencyjnie:
\begin{itemize}
\item każda zmienna $p\in\mathrm{Var}$ i stała $\bot$
      są formułami;
\item jeżeli $A,B$ są formułami, to
      $(A\wedge B)$, $(A\vee B)$ i $(A\to B)$
      są formułami;
\item jeżeli $A$ jest formułą, to $\Box A$ jest formułą.
\end{itemize}
Stosujemy skróty
\[
 \neg A=A\to\bot,\qquad
 \top=\neg\bot,\qquad
 A\leftrightarrow B=(A\to B)\wedge(B\to A),\qquad
 \Diamond A=\neg\Box\neg A.
\]
Operator $\Box$ nazywamy operatorem konieczności,
a $\Diamond$ operatorem możliwości.
\end{definicja}
Spójniki zdaniowe są w tym rozdziale klasyczne.
W szczególności nie używamy intuicjonistycznej
negacji do definiowania operatora możliwości.
Wyrażenia $\Box A$ oraz $\Diamond A$ nie są
nowymi zmiennymi, lecz formułami utworzonymi
z $A$.

\begin{definicja}[Głębokość modalna]
\label{def:ncl-modal-depth}
Głębokość modalną $\operatorname{md}(A)$ formuły
określamy przez
\[
\begin{aligned}
 \operatorname{md}(p)&=\operatorname{md}(\bot)=0,\\
 \operatorname{md}(A\circ B)
 &=\max(\operatorname{md}(A),\operatorname{md}(B)),
 \qquad \circ\in\{\wedge,\vee,\to\},\\
 \operatorname{md}(\Box A)&=1+\operatorname{md}(A).
\end{aligned}
\]
W szczególności $\operatorname{md}(\Diamond A)
=1+\operatorname{md}(A)$, co wynika z definicji
skrótu $\Diamond$.
\end{definicja}

\begin{przyklad}[Zakres operatorów modalnych]
\label{ex:ncl-modal-scope}
Formuła $\Box(p\to q)$ mówi o konieczności
implikacji. Formuła $\Box p\to\Box q$ mówi,
że konieczność $p$ pociąga konieczność $q$.
Obie mają głębokość modalną $1$, ale nie są
tym samym wyrażeniem.
Natomiast $\Box\Diamond p$ ma głębokość $2$:
wyraża konieczność możliwości $p$.
Nie należy jej utożsamiać z $\Diamond\Box p$,
czyli możliwością konieczności $p$.
\end{przyklad}

\begin{definicja}[Podstawienie modalne]
\label{def:ncl-modal-substitution}
Podstawienie $\sigma:\mathrm{Var}\to\mathrm{Fm}_{\Box}$
rozszerzamy na formuły, zachowując stałą $\bot$,
wszystkie spójniki dwuargumentowe i operator $\Box$.
W szczególności
\[
 \sigma(\Box A)=\Box\sigma(A),\qquad
 \sigma(\Diamond A)=\Diamond\sigma(A).
\]
Druga równość wynika z pierwszej i z zachowania
spójników występujących w definicji $\Diamond$.
\end{definicja}
Podstawienie jest jednoczesne: formuł wstawionych
za zmienne nie podstawiamy ponownie w tym samym
kroku. Dla podstawień $\sigma,\tau$ ich złożenie
na zmiennych określa wzór
$(\tau\circ\sigma)(p)=\tau(\sigma(p))$.
Indukcja po budowie formuły pokazuje, że
$(\tau\circ\sigma)(A)=\tau(\sigma(A))$:
dla zmiennych jest to definicja, dla stałej
obie strony są $\bot$, a dla spójników i $\Box$
korzystamy z hipotezy indukcyjnej dla argumentów.

\section{Normalna logika modalna K}

\begin{definicja}[System Hilberta K]
\label{def:ncl-modal-K}
System $\mathrm K$ ma następujące aksjomaty i reguły:
\begin{enumerate}
\item Wszystkie instancje klasycznych tautologii
      zdaniowych przez podstawienia formuł modalnych.
\item Schemat dystrybucji konieczności:
\[
 \tag{K}\label{eq:ncl-modal-axiom-K}
 \Box(A\to B)\to(\Box A\to\Box B).
\]
\item Reguła modus ponens:
\[
 \frac{A\qquad A\to B}{B}.
\]
\item Reguła konieczności (necessitation):
\[
 \frac{A}{\Box A}\quad(\mathrm{Nec}),
\]
      stosowana do twierdzeń, czyli formuł
      dowiedzionych bez dodatkowych przesłanek.
\end{enumerate}
Dowód twierdzenia jest skończonym ciągiem formuł,
z których każda jest aksjomatem lub wynika
z wcześniejszych wierszy przez jedną z reguł.
Piszemy $\vdash_{\mathrm K}A$, gdy istnieje taki
dowód o ostatnim wierszu $A$.
\end{definicja}
Schemat K pozwala przenosić konieczną implikację
na konieczność jej członów. Nie daje jednak
odwrotnej implikacji. W szczególności z samego
$A\to B$ jako lokalnego założenia nie wolno
przejść do $\Box(A\to B)$.

\begin{definicja}[Normalna logika modalna]
\label{def:ncl-modal-normal-logic}
Normalną logiką modalną nazywamy zbiór $L$
formuł modalnych zawierający klasyczne instancje
tautologii oraz wszystkie instancje K i domknięty
względem modus ponens, reguły konieczności
oraz dowolnych podstawień.
Piszemy $\vdash_L A$, gdy $A\in L$.
\end{definicja}

\begin{stwierdzenie}[Minimalność systemu K]
\label{prop:ncl-modal-K-minimal}
Zbiór twierdzeń systemu $\mathrm K$ jest najmniejszą
normalną logiką modalną.
\end{stwierdzenie}
\begin{proof}
Aksjomaty systemu są domknięte względem podstawień.
Podstawiając formuły we wszystkich wierszach dowodu,
otrzymujemy ponownie dowód: podstawienie zachowuje
zarówno modus ponens, jak i przejście od $A$ do
$\Box A$. Stąd zbiór twierdzeń K jest domknięty
względem podstawień. Domknięcie względem dwóch
pozostałych reguł wynika z możliwości połączenia
dowodów i dopisania odpowiedniego wniosku.
Jeżeli $L$ jest normalna, to indukcja po długości
dowodu K pokazuje, że każdy jego wiersz należy
do $L$. Zatem każdy normalny system zawiera
wszystkie twierdzenia K.
\end{proof}
Dla dowolnej rodziny dodatkowych schematów
aksjomatycznych piszemy $\mathrm K+\Sigma$ na
oznaczenie najmniejszej normalnej logiki
zawierającej te schematy. Istnieje ona jako
przecięcie wszystkich normalnych logik
zawierających $\Sigma$; rodzina ta jest niepusta,
bo zawiera zbiór wszystkich formuł.
Równoważnie bierzemy system Hilberta K
z dodatkowymi instancjami schematów: powtarzając
dowód minimalności otrzymujemy tę samą logikę.

\section{Reguły pochodne i rachunek operatorów}

\begin{twierdzenie}[Monotoniczność i zastępowanie równoważników]
\label{thm:ncl-modal-monotonicity}
W każdej normalnej logice modalnej $L$ zachodzą
reguły pochodne
\[
\begin{aligned}
 \vdash_L A\to B&\Longrightarrow
                   \vdash_L\Box A\to\Box B,\\
 \vdash_L A\to B&\Longrightarrow
                   \vdash_L\Diamond A\to\Diamond B,\\
 \vdash_L A\leftrightarrow B&\Longrightarrow
                   \vdash_L\Box A\leftrightarrow\Box B,\\
 \vdash_L A\leftrightarrow B&\Longrightarrow
                   \vdash_L\Diamond A\leftrightarrow\Diamond B.
\end{aligned}
\]
\end{twierdzenie}
\begin{proof}
Z twierdzenia $A\to B$ przez konieczność otrzymujemy
$\Box(A\to B)$, a następnie przez K i modus ponens
$\Box A\to\Box B$.
Klasyczna kontrapozycja daje z $A\to B$
twierdzenie $\neg B\to\neg A$.
Pierwsza reguła zastosowana do niego daje
$\Box\neg B\to\Box\neg A$.
Ponowna kontrapozycja daje
$\neg\Box\neg A\to\neg\Box\neg B$, czyli
$\Diamond A\to\Diamond B$.
Dwie ostatnie reguły wynikają z zastosowania
dwóch pierwszych do obu kierunków równoważności
i połączenia otrzymanych implikacji koniunkcją,
co jest dozwolone przez klasyczne tautologie
oraz modus ponens.
\end{proof}

\begin{wniosek}[Zastępowanie w kontekstach modalnych]
\label{cor:ncl-modal-context-replacement}
Niech $C[-]$ oznacza formułę z wyróżnionym miejscem
na podformułę. Jeśli $\vdash_L A\leftrightarrow B$,
to $\vdash_L C[A]\leftrightarrow C[B]$.
\end{wniosek}
\begin{proof}
Indukujemy po budowie kontekstu.
Dla samego wyróżnionego miejsca otrzymujemy
założenie. Dla spójników klasycznych korzystamy
z tautologii wyrażających zachowanie równoważności
po zastąpieniu jednego argumentu.
Dla kontekstu $\Box D[-]$ stosujemy hipotezę
indukcyjną i trzecią regułę z poprzedniego
twierdzenia. Skróty, w tym $\Diamond$, są objęte
indukcją po rozwinięciu.
\end{proof}

\begin{twierdzenie}[Prawa normalnych operatorów modalnych]
\label{thm:ncl-modal-normal-laws}
W każdej normalnej logice modalnej dowodliwe są
formuły
\[
\begin{gathered}
 \Box\top,\qquad \neg\Diamond\bot,\\
 \Box(A\wedge B)\leftrightarrow(\Box A\wedge\Box B),\\
 \Diamond(A\vee B)\leftrightarrow
                  (\Diamond A\vee\Diamond B),\\
 (\Box A\vee\Box B)\to\Box(A\vee B),\\
 \Diamond(A\wedge B)\to(\Diamond A\wedge\Diamond B),\\
 \neg\Box A\leftrightarrow\Diamond\neg A,\qquad
 \neg\Diamond A\leftrightarrow\Box\neg A.
\end{gathered}
\]
\end{twierdzenie}
\begin{proof}
Z tautologii $\top$ przez regułę konieczności
otrzymujemy $\Box\top$.
Ponieważ $\neg\bot=\top$, klasyczna eliminacja
podwójnej negacji daje $\neg\Diamond\bot$.

Z tautologii $A\wedge B\to A$ oraz
$A\wedge B\to B$ przez monotoniczność $\Box$
otrzymujemy obie implikacje potrzebne do
$\Box(A\wedge B)\to(\Box A\wedge\Box B)$.
W drugą stronę stosujemy konieczność do tautologii
$A\to(B\to A\wedge B)$, a następnie dwukrotnie K.
Otrzymujemy kolejno
\[
 \Box A\to\Box(B\to A\wedge B),\qquad
 \Box(B\to A\wedge B)\to
       (\Box B\to\Box(A\wedge B)).
\]
Składanie implikacji i klasyczne przekształcenie
$X\to(Y\to Z)$ w $(X\wedge Y)\to Z$ dają
$(\Box A\wedge\Box B)\to\Box(A\wedge B)$.
Równoważność dla możliwości sumy wynika z definicji
$\Diamond$, praw de Morgana, zastępowania
równoważników oraz właśnie udowodnionej
równoważności dla $\Box$ i koniunkcji.

Tautologie $A\to A\vee B$ i $B\to A\vee B$
przenosimy przez $\Box$ i łączymy klasyczną
eliminacją alternatywy. Analogicznie z implikacji
$A\wedge B\to A$ i $A\wedge B\to B$ przez
monotoniczność $\Diamond$ otrzymujemy implikację
dla możliwości koniunkcji.
Ostatnie dwie równoważności wynikają z definicji
$\Diamond$, klasycznej podwójnej negacji i
zastępowania równoważników pod $\Box$.
\end{proof}
Indukcja po liczbie członów rozszerza prawo
konieczności koniunkcji na dowolną skończoną
rodzinę formuł, a prawo możliwości alternatywy
na dowolną skończoną alternatywę.
Przypadki pustych rodzin to odpowiednio
$\Box\top$ i $\neg\Diamond\bot$.

\begin{stwierdzenie}[Modalny modus ponens]
\label{prop:ncl-modal-mixed-mp}
W każdej normalnej logice modalnej zachodzi
\[
 \vdash_L\Box(A\to B)\to(\Diamond A\to\Diamond B).
\]
\end{stwierdzenie}
\begin{proof}
Z tautologii $(A\to B)\to(\neg B\to\neg A)$
i monotoniczności $\Box$ otrzymujemy
$\Box(A\to B)\to\Box(\neg B\to\neg A)$.
Schemat K daje
$\Box(\neg B\to\neg A)\to
(\Box\neg B\to\Box\neg A)$.
Składamy obie implikacje i stosujemy klasyczną
kontrapozycję do implikacji w następniku.
Wynik to
$\Box(A\to B)\to
(\neg\Box\neg A\to\neg\Box\neg B)$.
\end{proof}

\section{Podstawowe rozszerzenia aksjomatyczne}

\begin{definicja}[Schematy D, T, B, 4 i 5]
\label{def:ncl-modal-axiom-schemes}
Rozważamy następujące schematy, w których $A$
jest dowolną formułą modalną:
\[
\begin{array}{c|l}
 \mathrm D&\Box A\to\Diamond A\\
 \mathrm T&\Box A\to A\\
 \mathrm B&A\to\Box\Diamond A\\
 4&\Box A\to\Box\Box A\\
 5&\Diamond A\to\Box\Diamond A.
\end{array}
\]
Definiujemy logiki
\[
\begin{aligned}
 \mathrm{KD}&=\mathrm K+\mathrm D,&
 \mathrm T&=\mathrm K+\mathrm T,\\
 \mathrm{KB}&=\mathrm K+\mathrm B,&
 \mathrm{K4}&=\mathrm K+4,\\
 \mathrm{K5}&=\mathrm K+5,&
 \mathrm B&=\mathrm K+\mathrm T+\mathrm B,\\
 \mathrm{S4}&=\mathrm K+\mathrm T+4,&
 \mathrm{S5}&=\mathrm K+\mathrm T+5.
\end{aligned}
\]
Po prawej stronie symbole D, T, B, 4, 5 oznaczają
schematy; po lewej nazwy systemów. Oznaczenia
$\mathrm{KB}$ i $\mathrm B$ rozróżniamy: drugi
system ma dodatkowo schemat T.
\end{definicja}

Schemat D wyklucza konieczność zdania wraz
z niemożliwością tego zdania.
Schemat T wyraża faktywność konieczności.
Schemat B pozwala z prawdziwości zdania przejść
do konieczności jego możliwości.
Schemat 4 wyraża przejście od konieczności
do konieczności konieczności, a schemat 5
od możliwości do konieczności możliwości.
Są to intuicje syntaktyczne; ich dokładne znaczenie
w kategoriach własności relacji dostępności
zostanie udowodnione w rozdziale o semantyce Kripkego.

\begin{stwierdzenie}[Dualne postacie schematów]
\label{prop:ncl-modal-dual-axioms}
W obecności K następujące pary schematów
wyznaczają te same rozszerzenia:
\[
\begin{array}{c|l}
 \mathrm D & \neg\Box\bot\\
 \mathrm T & A\to\Diamond A\\
 \mathrm B & \Diamond\Box A\to A\\
 4 & \Diamond\Diamond A\to\Diamond A\\
 5 & \Diamond\Box A\to\Box A .
\end{array}
\]
\end{stwierdzenie}
\begin{proof}
D dla $A=\bot$ daje $\Box\bot\to\Diamond\bot$,
a ponieważ $\neg\Diamond\bot$ jest twierdzeniem K,
otrzymujemy $\neg\Box\bot$.
Odwrotnie, z prawa konieczności koniunkcji i
zastępowania równoważników mamy
$(\Box A\wedge\Box\neg A)\to\Box\bot$.
Z $\neg\Box\bot$ wynika więc
$\Box A\to\neg\Box\neg A$, czyli D.
Dla T podstawiamy $\neg A$ do schematu
$\Box A\to A$, stosujemy kontrapozycję i
eliminujemy podwójną negację, otrzymując
$A\to\Diamond A$. Te same operacje zastosowane
do postaci dualnej dają T.
Dla B z instancji
$\neg A\to\Box\Diamond\neg A$ przez
kontrapozycję otrzymujemy
$\neg\Box\Diamond\neg A\to A$ po uproszczeniu
podwójnej negacji. Z definicji i praw dualności
$\neg\Box\Diamond\neg A$ jest równoważne
$\Diamond\Box A$. Odwrócenie argumentu
daje pierwotny schemat B.
Dla 4 kontrapozycja instancji
$\Box\neg A\to\Box\Box\neg A$ daje
$\Diamond\Diamond A\to\Diamond A$
po zastosowaniu praw dualności.
Dla 5 analogicznie kontrapozycja instancji
$\Diamond\neg A\to\Box\Diamond\neg A$
daje $\Diamond\Box A\to\Box A$.
W dwóch ostatnich przypadkach ponowne
podstawienie negacji i kontrapozycja dają
schematy wyjściowe, ponieważ prawa dualności
oraz podwójnej negacji są dowodliwe już w K.
\end{proof}

\begin{wniosek}[Faktywność pociąga D]
\label{cor:ncl-modal-T-implies-D}
Każda normalna logika zawierająca schemat T
zawiera także D.
\end{wniosek}
\begin{proof}
Z T mamy $\Box A\to A$, a z postaci dualnej T
ze Stwierdzenia~\ref{prop:ncl-modal-dual-axioms}
mamy $A\to\Diamond A$. Składamy implikacje,
korzystając z klasycznej tautologii przechodniości
implikacji i modus ponens.
\end{proof}

\section{Konsekwencja lokalna i globalna}

\begin{definicja}[Dwie relacje konsekwencji syntaktycznej]
\label{def:ncl-modal-local-global}
Niech $L$ będzie normalną logiką modalną, a
$\Gamma\subseteq\mathrm{Fm}_{\Box}$.
Piszemy $\Gamma\vdash_L^{\mathrm{lok}} A$, jeśli
istnieje skończony ciąg kończący się formułą $A$,
którego każdy wiersz jest elementem $\Gamma$,
twierdzeniem $L$ albo wynika z wcześniejszych
wierszy przez modus ponens.
Piszemy $\Gamma\vdash_L^{\mathrm{glob}} A$,
jeśli dodatkowo wolno stosować regułę konieczności
do każdego wcześniejszego wiersza, również do
formuł uzyskanych z przesłanek.
\end{definicja}
W obu definicjach podstawienie nie jest regułą
stosowaną do dowolnej przesłanki. Z faktu, że
$p$ jest założeniem, nie wynika dowolna formuła
$\sigma(p)$. Domknięcie logiki względem podstawień
dotyczy zbioru twierdzeń $L$; własność strukturalna
konsekwencji polega natomiast na jednoczesnym
podstawieniu w przesłankach i wniosku.

W konsekwencji lokalnej zakładamy prawdziwość
przesłanek w aktualnym stanie. Konsekwencja
globalna odpowiada założeniu ich prawdziwości
we wszystkich stanach modelu. Jest to motywacja
dla dwóch różnych rachunków syntaktycznych;
ich semantyczną pełność rozważymy osobno.

\begin{stwierdzenie}[Własności strukturalne obu konsekwencji]
\label{prop:ncl-modal-consequence-structural}
Dla każdej z relacji $\vdash_L^{\mathrm{lok}}$,
$\vdash_L^{\mathrm{glob}}$ zachodzą:
\begin{enumerate}
\item zwrotność: $A\in\Gamma$ implikuje
      $\Gamma\vdash A$;
\item monotoniczność: z $\Gamma\subseteq\Delta$
      i $\Gamma\vdash A$ wynika $\Delta\vdash A$;
\item cięcie: jeżeli $\Gamma\vdash B$ dla każdego
      $B\in\Delta$ oraz $\Delta\vdash A$, to
      $\Gamma\vdash A$;
\item skończoność: jeżeli $\Gamma\vdash A$,
      to pewien skończony $\Gamma_0\subseteq\Gamma$
      spełnia $\Gamma_0\vdash A$;
\item strukturalność: $\Gamma\vdash A$ implikuje
      $\sigma(\Gamma)\vdash\sigma(A)$.
\end{enumerate}
Ponadto
\[
 \varnothing\vdash_L^{\mathrm{lok}} A
 \iff \varnothing\vdash_L^{\mathrm{glob}} A
 \iff A\in L.
\]
\end{stwierdzenie}
\begin{proof}
Zwrotność i monotoniczność wynikają wprost
z definicji dowodu. Dla cięcia bierzemy dowód
$A$ z $\Delta$ i dla każdej użytej w nim przesłanki
$B\in\Delta$ wybieramy jej dowód z $\Gamma$.
Dołączamy te skończenie liczne dowody przed
dowodem $A$. Każdy wiersz będący przesłanką
z $\Delta$ ma w ten sposób wcześniejsze
uzasadnienie z $\Gamma$, a pozostałe wiersze
można odtworzyć przez dozwolone reguły.
Jeśli ostatni wiersz jest już wcześniej uzyskany,
możemy zakończyć dowód na jego wystąpieniu.
W przypadku globalnym wszystkie zastosowania
konieczności pozostają legalne, ponieważ
dopuszczamy tę regułę także dla formuł
uzyskanych z $\Gamma$.
Skończoność wynika ze skończoności dowodu.
Podstawiając jednocześnie we wszystkich wierszach,
zachowujemy przynależność do zbioru twierdzeń $L$,
modus ponens i, w przypadku globalnym, regułę
konieczności.
Dla pustego zbioru przesłanek indukcja po
wierszach dowodu pokazuje, że każdy wiersz
należy do $L$, ponieważ $L$ jest domknięta
względem obu reguł. Odwrotnie każde twierdzenie
$L$ jest jednoelementowym dowodem w obu rachunkach.
\end{proof}

\begin{twierdzenie}[Lokalne twierdzenie o dedukcji]
\label{thm:ncl-modal-local-deduction}
Dla dowolnego zbioru $\Gamma$ i formuł $A,B$
zachodzi
\[
 \Gamma\cup\{A\}\vdash_L^{\mathrm{lok}} B
 \iff \Gamma\vdash_L^{\mathrm{lok}} A\to B.
\]
\end{twierdzenie}
\begin{proof}
Kierunek od prawej do lewej wynika
z monotoniczności oraz modus ponens.
W drugim kierunku indukujemy po wierszach $C$
lokalnego dowodu $B$ z $\Gamma\cup\{A\}$.
Jeśli $C=A$, to $A\to C$ jest klasyczną
tautologią.
Jeśli $C\in\Gamma$ albo $C\in L$, to
z $C$ i tautologii $C\to(A\to C)$ uzyskujemy
$A\to C$ z $\Gamma$.
Jeżeli $C$ wynika z wcześniejszych $D,D\to C$,
to mamy z hipotezy indukcyjnej dowody
$A\to D$ oraz $A\to(D\to C)$ z $\Gamma$.
Stosujemy klasyczną tautologię
\[
 (A\to(D\to C))\to((A\to D)\to(A\to C))
\]
oraz dwukrotnie modus ponens.
Nie występuje przypadek konieczności
stosowanej do przesłanki: zastosowania
konieczności w dowodach twierdzeń $L$
są już zawarte w przypadku $C\in L$.
\end{proof}

\begin{wniosek}[Skończone przesłanki lokalne]
\label{cor:ncl-modal-local-finite}
Dla skończonego zbioru $\Delta$ i formuły $B$
zachodzi
\[
 \Delta\vdash_L^{\mathrm{lok}} B
 \iff
 \vdash_L \left(\bigwedge_{A\in\Delta} A\right)\to B.
\]
Pusta koniunkcja jest równa $\top$.
\end{wniosek}
\begin{proof}
Dla niepustej $\Delta$ możemy zastąpić wszystkie
jej elementy ich koniunkcją: każdy składnik
wynika z koniunkcji przez klasyczną tautologię
i modus ponens, a koniunkcja wynika z jej
składników przez iterowane wprowadzanie
koniunkcji, które jest regułą pochodną rachunku
klasycznego. Własność cięcia uzasadnia oba
kierunki zamiany.
Stosujemy Twierdzenie~\ref{thm:ncl-modal-local-deduction}
dla $\Gamma=\varnothing$ oraz
Stwierdzenie~\ref{prop:ncl-modal-consequence-structural}.
Dla $\Delta=\varnothing$ korzystamy z dowodliwości
$\top$ i równoważności $B\leftrightarrow(\top\to B)$.
\end{proof}

\subsection{Iteracje konieczności i dedukcja globalna}

\begin{definicja}[Domknięcie przesłanek względem konieczności]
\label{def:ncl-modal-box-closure}
Piszemy $\Box^0 A=A$ oraz
$\Box^{n+1}A=\Box(\Box^n A)$ dla $n\geq0$.
Dla zbioru formuł $\Gamma$ określamy
\[
 \Gamma^\Box=\{\Box^n A:A\in\Gamma,\ n\geq0\},
 \qquad
 A^{[n]}=\bigwedge_{k=0}^n\Box^k A.
\]
\end{definicja}

\begin{twierdzenie}[Redukcja konsekwencji globalnej do lokalnej]
\label{thm:ncl-modal-global-local}
Dla dowolnej normalnej logiki $L$ zachodzi
\[
 \Gamma\vdash_L^{\mathrm{glob}} B
 \iff
 \Gamma^\Box\vdash_L^{\mathrm{lok}} B.
\]
\end{twierdzenie}
\begin{proof}
Każda formuła $\Box^n A$ z $\Gamma^\Box$
jest globalnie dowodliwa z $\Gamma$ przez
$n$ zastosowań konieczności.
Lokalny dowód z $\Gamma^\Box$ jest również
dowodem globalnym z tego zbioru, więc przez
cięcie otrzymujemy jeden kierunek.

Dla odwrotnego kierunku zauważmy, że jeżeli
$\Gamma^\Box\vdash_L^{\mathrm{lok}} C$, to
$\Gamma^\Box\vdash_L^{\mathrm{lok}}\Box C$.
Istotnie, lokalny dowód $C$ wykorzystuje
skończenie wiele przesłanek $D_1,\ldots,D_m$
z $\Gamma^\Box$. Wniosek~\ref{cor:ncl-modal-local-finite}
daje
\[
 \vdash_L(D_1\wedge\cdots\wedge D_m)\to C.
\]
Przez monotoniczność konieczności otrzymujemy
\[
 \vdash_L\Box(D_1\wedge\cdots\wedge D_m)\to\Box C.
\]
Każda $\Box D_i$ należy do $\Gamma^\Box$.
Prawo konieczności skończonej koniunkcji
(Twierdzenie~\ref{thm:ncl-modal-normal-laws})
pozwala z nich lokalnie uzyskać
$\Box(D_1\wedge\cdots\wedge D_m)$, a następnie
$\Box C$. Dla $m=0$ używamy twierdzenia
$\Box\top$.
Zatem lokalne konsekwencje $\Gamma^\Box$
są domknięte względem konieczności.
Indukcja po dowodzie globalnym z $\Gamma$
pokazuje teraz, że każdy jego wiersz jest
lokalnie dowodliwy z $\Gamma^\Box$:
przesłanki należą do $\Gamma^\Box$, twierdzenia
i modus ponens są dozwolone lokalnie, a przypadek
konieczności obsługuje właśnie wykazana własność.
\end{proof}

\begin{twierdzenie}[Globalne twierdzenie o dedukcji]
\label{thm:ncl-modal-global-deduction}
Dla dowolnego zbioru $\Gamma$ i formuł $A,B$
zachodzi
\[
 \Gamma\cup\{A\}\vdash_L^{\mathrm{glob}} B
 \iff
 \exists n\geq0\quad
 \Gamma\vdash_L^{\mathrm{glob}} A^{[n]}\to B.
\]
Równoważnie po prawej stronie można przyjąć
$\Gamma^\Box\vdash_L^{\mathrm{lok}} A^{[n]}\to B$.
\end{twierdzenie}
\begin{proof}
Z Twierdzenia~\ref{thm:ncl-modal-global-local}
lewa strona jest równoważna lokalnej dowodliwości
$B$ z $\Gamma^\Box\cup\{\Box^k A:k\geq0\}$.
Dowód wykorzystuje skończenie wiele iteracji
$\Box^k A$. Wybieramy $n$ nie mniejsze od
wszystkich ich wykładników (jeśli nie użyto
żadnej, bierzemy $n=0$).
Każda z wykorzystanych iteracji jest lokalną
konsekwencją $A^{[n]}$ na mocy klasycznych
tautologii eliminacji koniunkcji.
Cięcie daje
$\Gamma^\Box\cup\{A^{[n]}\}\vdash_L^{\mathrm{lok}}B$.
Lokalne twierdzenie o dedukcji i redukcja
konsekwencji globalnej do lokalnej dają
obie wersje prawej strony.

Odwrotnie, jeśli prawa strona zachodzi, to
globalnie z $\Gamma\cup\{A\}$ dowodzimy
każdej $\Box^k A$ dla $0\leq k\leq n$,
a następnie ich koniunkcji. Z monotoniczności
mamy również dowód $A^{[n]}\to B$ z tych
samych przesłanek. Modus ponens daje $B$.
\end{proof}

\begin{wniosek}[Dedukcja globalna w K4 i S4]
\label{cor:ncl-modal-global-deduction-K4}
Jeżeli $L$ zawiera schemat 4, to
\[
 \Gamma\cup\{A\}\vdash_L^{\mathrm{glob}} B
 \iff
 \Gamma\vdash_L^{\mathrm{glob}}(A\wedge\Box A)\to B.
\]
Jeżeli ponadto $L$ zawiera T, to równoważnie
$\Gamma\vdash_L^{\mathrm{glob}}\Box A\to B$.
\end{wniosek}
\begin{proof}
Dla każdego $k\geq1$ mamy
$\vdash_L\Box A\to\Box^k A$.
Dowodzimy indukcyjnie: dla $k=1$ jest to
tautologia, a przejście od $k$ do $k+1$
wynika ze złożenia dotychczasowej implikacji
z instancją 4 dla $\Box^{k-1}A$.
Zatem $A\wedge\Box A$ pociąga każdą
skończoną koniunkcję $A^{[n]}$, również
dla $n=0$. Stosujemy
Twierdzenie~\ref{thm:ncl-modal-global-deduction}.
Odwrotny kierunek otrzymujemy z tego
samego twierdzenia dla $n=1$.
Przy T formuły $A\wedge\Box A$ i $\Box A$
są dowodliwie równoważne.
\end{proof}

\subsection{Dlaczego rozróżnienie jest konieczne}

Do kontrprzykładów wystarczy elementarna
interpretacja relacyjna. Niech $W$ będzie
niepustym zbiorem stanów, $R\subseteq W\times W$,
a zmienne przyjmują dowolne podzbiory $W$.
Spójniki klasyczne interpretujemy punktowo,
a $\Box C$ jest prawdziwa w stanie $w$,
gdy $C$ jest prawdziwa we wszystkich stanach
$z$ spełniających $wRz$.
Schemat K jest prawdziwy w każdym stanie:
jeśli wszystkie dostępne stany spełniają
$A\to B$ oraz wszystkie spełniają $A$,
to wszystkie spełniają $B$.
Tautologie są prawdziwe punktowo, a modus
ponens zachowuje prawdziwość w danym stanie.
Konieczność zachowuje prawdziwość we wszystkich
stanach: gdy $C$ jest wszędzie prawdziwa, to
$\Box C$ również. Indukcja po dowodach K daje
zatem prawdziwość wszystkich twierdzeń K
we wszystkich stanach każdego takiego modelu.
W szczególności lokalny dowód z przesłanek
zachowuje prawdziwość w stanie, w którym
prawdziwe są wszystkie użyte przesłanki.
Pełny rozwój tej semantyki pozostawiamy
następnemu rozdziałowi.

\begin{przyklad}[Konieczność nie jest lokalną regułą wnioskowania]
\label{ex:ncl-modal-local-global-difference}
Mamy $p\vdash_{\mathrm K}^{\mathrm{glob}}\Box p$
przez jedno zastosowanie konieczności, lecz
$p\nvdash_{\mathrm K}^{\mathrm{lok}}\Box p$.
Rozważmy dwa stany $w,z$ z jedyną parą
$wRz$ i niech $p$ będzie prawdziwe tylko w $w$.
W stanie $w$ przesłanka jest prawdziwa,
ale $\Box p$ jest fałszywa.
W szczególności zwykłe twierdzenie o dedukcji
nie zachodzi dla konsekwencji globalnej:
$p\to\Box p$ nie jest twierdzeniem K,
chociaż $p\vdash_{\mathrm K}^{\mathrm{glob}}\Box p$.
\end{przyklad}

\begin{center}
\begin{tikzpicture}[>=Stealth]
\node[draw,circle,minimum size=10mm,fill=blue!10] (w) at (0,0) {$w$};
\node[draw,circle,minimum size=10mm] (z) at (3,0) {$z$};
\draw[->,thick] (w)--node[above] {$R$}(z);
\node[below=4pt of w] {$p,\ \neg\Box p$};
\node[below=4pt of z] {$\neg p$};
\end{tikzpicture}
\end{center}

\begin{przyklad}[Brak jednolitego ograniczenia liczby iteracji w K]
\label{ex:ncl-modal-global-no-uniform-bound}
Dla każdego $n\geq0$ zachodzi
$p\vdash_{\mathrm K}^{\mathrm{glob}}\Box^{n+1}p$,
ale
\[
 \nvdash_{\mathrm K} p^{[n]}\to\Box^{n+1}p.
\]
Weźmy stany $0,1,\ldots,n+1$ z relacją
$iRj$ wtedy i tylko wtedy, gdy $j=i+1$.
Niech $p$ będzie prawdziwe w stanach
$0,\ldots,n$, ale fałszywe w stanie $n+1$.
Dla $0\leq k\leq n+1$ istnieje dokładnie jedna
ścieżka długości $k$ z $0$, kończąca się w $k$.
Indukcja po długości ścieżki, stosująca
klauzulę $\Box$, daje, że $\Box^k p$ jest
prawdziwa w $0$ dokładnie wtedy, gdy $p$
jest prawdziwe w $k$.
Zatem w $0$ prawdziwa jest cała koniunkcja
$p^{[n]}$, ale fałszywa jest $\Box^{n+1}p$.
Poprawność twierdzeń K w powyższej interpretacji
wyklucza dowodliwość implikacji.
Oznacza to, że liczba iteracji w globalnym
twierdzeniu o dedukcji nie może być z góry
ograniczona jedną liczbą niezależną od wniosku.
\end{przyklad}

\begin{wniosek}[Domknięcie teorii lokalnej i globalnej]
\label{cor:ncl-modal-theory-closure}
Niech $\operatorname{Cn}^{\mathrm{lok}}_L(\Gamma)$
oraz $\operatorname{Cn}^{\mathrm{glob}}_L(\Gamma)$
oznaczają zbiory odpowiednich konsekwencji.
Oba operatory są ekstensywne, monotoniczne,
idempotentne i finitarne. Ponadto
\[
 \operatorname{Cn}^{\mathrm{glob}}_L(\Gamma)
 =\operatorname{Cn}^{\mathrm{lok}}_L(\Gamma^\Box).
\]
Zbiór domknięty względem konsekwencji lokalnej
jest domknięty względem konsekwencji globalnej
dokładnie wtedy, gdy jest domknięty względem
operatora $\Box$.
\end{wniosek}
\begin{proof}
Pierwsze cztery własności wynikają kolejno
ze zwrotności, monotoniczności, cięcia wraz
ze zwrotnością oraz skończoności
(Stwierdzenie~\ref{prop:ncl-modal-consequence-structural}).
Równość operatorów jest treścią
Twierdzenia~\ref{thm:ncl-modal-global-local}.
Jeżeli $T$ jest domknięty lokalnie i
$A\in T$ pociąga $\Box A\in T$, to
$T^\Box=T$, więc równość operatorów daje
domknięcie globalne.
Odwrotnie domknięcie globalne daje domknięcie
względem $\Box$ bezpośrednio z reguły konieczności.
\end{proof}

\section{Hierarchia normalnych logik modalnych}

Porównujemy logiki przez inkluzję ich zbiorów twierdzeń.
Jeżeli $L_1\subseteq L_2$, każda formuła dowodliwa
w $L_1$ jest dowodliwa również w $L_2$.
Nie jest to porównanie liczby aksjomatów:
różne zbiory aksjomatów mogą generować tę samą
normalną logikę.

\begin{lemat}[Zależności między B, T, 4 i 5]
\label{lem:ncl-modal-axiom-dependencies}
W normalnej logice modalnej zachodzą następujące
zależności między schematami:
\[
\begin{aligned}
 \mathrm T+5&\Longrightarrow\mathrm B,\\
 \mathrm B+4&\Longrightarrow5,\\
 \mathrm B+5&\Longrightarrow4.
\end{aligned}
\]
Każda strzałka oznacza dowodliwość wszystkich
instancji schematu po prawej stronie przy założeniu
schematów po lewej.
\end{lemat}
\begin{proof}
Dla pierwszej zależności bierzemy dualną postać T
ze Stwierdzenia~\ref{prop:ncl-modal-dual-axioms}
i schemat 5:
\[
 A\to\Diamond A,\qquad
 \Diamond A\to\Box\Diamond A.
\]
Złożenie implikacji daje $A\to\Box\Diamond A$,
czyli B.

Dla drugiej stosujemy B do $\Diamond A$:
\[
 \Diamond A\to\Box\Diamond\Diamond A.
\]
Dualna postać 4 to $\Diamond\Diamond A\to\Diamond A$.
Monotoniczność konieczności
(Twierdzenie~\ref{thm:ncl-modal-monotonicity})
daje
$\Box\Diamond\Diamond A\to\Box\Diamond A$.
Po złożeniu otrzymujemy 5.

Dla trzeciej stosujemy B do $\Box A$:
\[
 \Box A\to\Box\Diamond\Box A.
\]
Dualna postać 5 to $\Diamond\Box A\to\Box A$.
Przeniesienie jej przez konieczność daje
$\Box\Diamond\Box A\to\Box\Box A$.
Złożenie implikacji daje 4.
We wszystkich trzech przypadkach złożenie
jest uzasadnione klasyczną tautologią
przechodniości implikacji i modus ponens.
\end{proof}

\begin{twierdzenie}[Alternatywne aksjomatyzacje S5]
\label{thm:ncl-modal-S5-alternatives}
Następujące logiki są równe:
\[
 \mathrm{S5}
 =\mathrm K+\mathrm T+5
 =\mathrm K+\mathrm T+\mathrm B+4
 =\mathrm{S4}+5
 =\mathrm B+5
 =\mathrm K+\mathrm D+\mathrm B+4.
\]
\end{twierdzenie}
\begin{proof}
Pierwsza równość jest definicją S5.
Z T i 5 otrzymujemy B, a z B i 5 otrzymujemy 4
na mocy Lematu~\ref{lem:ncl-modal-axiom-dependencies}.
Odwrotnie B i 4 dają 5. Stąd równość
$\mathrm{S5}=\mathrm K+\mathrm T+\mathrm B+4$
i, korzystając z definicji systemów,
równości z $\mathrm{S4}+5$ oraz $\mathrm B+5$.

Pozostaje ostatnia aksjomatyzacja.
Wystarczy wyprowadzić T z D i B.
Z D dla $\Box A$ mamy
$\Box\Box A\to\Diamond\Box A$,
a z dualnej postaci B mamy
$\Diamond\Box A\to A$.
Z 4 otrzymujemy więc ciąg implikacji
\[
 \Box A\to\Box\Box A\to\Diamond\Box A\to A,
\]
który daje T w systemie $\mathrm K+\mathrm D+\mathrm B+4$.
W przeciwnym kierunku T daje D na mocy
Wniosku~\ref{cor:ncl-modal-T-implies-D}.
\end{proof}

\begin{uwaga}[D i B bez schematu 4]
\label{rem:ncl-modal-DB-not-T}
Same D i B nie wystarczają do wyprowadzenia T
w normalnej logice. Istotnie,
na ramie dwóch stanów połączonych w obie strony,
bez pętli, D i B są ważne, lecz T zawodzi.
Uzasadnienie ważności i kontrmodel podamy
w dalszej części sekcji.
Zatem w ostatniej aksjomatyzacji S5 w
Twierdzeniu~\ref{thm:ncl-modal-S5-alternatives}
schematu 4 nie można po prostu pominąć.
\end{uwaga}

\subsection{Skończone świadectwa ścisłości inkluzji}

Do sprawdzania niedowodliwości użyjemy
elementarnej interpretacji relacyjnej wprowadzonej
przy Przykładzie~\ref{ex:ncl-modal-local-global-difference}.
Nie rozwijamy tu jeszcze ogólnej teorii korespondencji
ani pełności. Potrzebujemy jedynie kierunku
od własności relacji do ważności odpowiedniego
schematu.

\begin{lemat}[Wystarczające warunki ważności schematów]
\label{lem:ncl-modal-frame-sufficient}
Niech $W$ będzie niepustym zbiorem stanów
i $R$ relacją na $W$. Dla każdego wartościowania
wszystkie instancje wskazanego schematu są
prawdziwe w każdym stanie, jeśli $R$ ma
odpowiednią własność:
\[
\begin{array}{c|l}
 \mathrm D&\forall w\,\exists z\quad wRz\\
 \mathrm T&\forall w\quad wRw\\
 \mathrm B& wRz\Longrightarrow zRw\\
 4&wRz,\ zRu\Longrightarrow wRu\\
 5&wRz,\ wRu\Longrightarrow zRu.
\end{array}
\]
\end{lemat}
\begin{proof}
Z klasycznej definicji możliwości wynika,
że $\Diamond A$ jest prawdziwa w $w$ dokładnie
wtedy, gdy istnieje $z$ z $wRz$, w którym
$A$ jest prawdziwa.

Dla D, jeżeli $\Box A$ jest prawdziwa w $w$,
wybieramy dostępny stan $z$. W nim $A$ jest
prawdziwa, więc w $w$ prawdziwa jest $\Diamond A$.
Dla T wystarczy wybrać $z=w$.
Dla B, jeżeli $A$ jest prawdziwa w $w$ i $wRz$,
to $zRw$, więc $\Diamond A$ jest prawdziwa w $z$.
Ponieważ $z$ było dowolnym dostępnym stanem,
$\Box\Diamond A$ jest prawdziwa w $w$.
Dla 4, jeśli $\Box A$ jest prawdziwa w $w$,
to dla dowolnych $wRz$ i $zRu$ zachodzi $wRu$,
a zatem $A$ jest prawdziwa w $u$.
W konsekwencji $\Box A$ jest prawdziwa w każdym
$z$ dostępnym z $w$.
Dla 5, jeżeli $\Diamond A$ jest prawdziwa w $w$,
wybieramy świadka $u$ z $wRu$, w którym $A$
jest prawdziwa. Dla dowolnego $z$ z $wRz$
mamy $zRu$, więc $\Diamond A$ jest prawdziwa w $z$.
To daje $\Box\Diamond A$ w $w$.

Poprawność K oraz zachowanie globalnej ważności
przez obie reguły dowodzenia sprawiają,
że na ramie spełniającej wymienione warunki
ważne są wszystkie twierdzenia odpowiedniego
rozszerzenia aksjomatycznego K, przez indukcję
po dowodzie. Nie zakładamy twierdzenia o pełności.
\end{proof}

\begin{twierdzenie}[Ścisłe inkluzje podstawowych systemów]
\label{thm:ncl-modal-hierarchy}
Zachodzą ścisłe inkluzje
\[
\begin{gathered}
 \mathrm K\subsetneq\mathrm{KD}\subsetneq\mathrm T,\\
 \mathrm T\subsetneq\mathrm B\subsetneq\mathrm{S5},
 \qquad
 \mathrm T\subsetneq\mathrm{S4}\subsetneq\mathrm{S5}.
\end{gathered}
\]
Systemy B i S4 są nieporównywalne.
Ponadto $\mathrm K\subsetneq\mathrm{K4}$,
$\mathrm K\subsetneq\mathrm{KB}$ i
$\mathrm K\subsetneq\mathrm{K5}$,
a systemy K4 i K5 są nieporównywalne.
\end{twierdzenie}
\begin{proof}
Inkluzje niestrikte wynikają z definicji,
Wniosku~\ref{cor:ncl-modal-T-implies-D}
oraz Twierdzenia~\ref{thm:ncl-modal-S5-alternatives}.
Dla ścisłości i nieporównywalności podamy
konkretne skończone ramy. W tabeli $R(i)$
oznacza zbiór stanów dostępnych z $i$;
stanami są dokładnie indeksy wymienionych wierszy
relacji. Zapis $V(p)$ określa zbiór stanów,
w których prawdziwe jest $p$. Kolumna $w$
wskazuje stan, w którym dana instancja
schematu dla $A=p$ jest fałszywa.
\[
\begin{array}{c|l|c|c|c|c}
 & R(i) & \text{ważne} & \text{zawodzi}
 &V(p)&w\\ \hline
 F_0&R(0)=\varnothing&B,4,5&D&\{0\}&0\\
 F_1&R(0)=R(1)=\{0\}&D,4,5&T&\{0\}&1\\
 F_2&R(0)=\{0,1\},\ R(1)=\{1\}&T,4,D&B&\{0\}&0\\
 F_3&
 \begin{array}{l}R(0)=\{0,1\}\\
 R(1)=\{0,1,2\}\\R(2)=\{1,2\}\end{array}
 &T,B,D&4&\{0,1\}&0\\
 F_4&R(0)=\{1\},\ R(1)=\{0\}&D,B&4&\{1\}&0\\
 F_5&R(0)=\{1\},\ R(1)=\varnothing&4&B,5&\{0,1\}&0\\
 F_6&
 \begin{array}{l}R(0)=R(1)=\{0,1\}\\R(2)=\{0\}\end{array}
 &D,5&4&\{0,2\}&2
\end{array}
\]
W kolumnie ważności wymieniono potrzebne
schematy, niekoniecznie wszystkie ważne na ramie.

Sprawdzamy świadectwa w tabeli.
Na $F_0$ mamy $\Box p$ prawdziwe próżniowo,
a $\Diamond p$ fałszywe.
Na $F_1$ stan $1$ widzi wyłącznie stan $0$,
więc $\Box p$ jest prawdziwe, a $p$ fałszywe
w stanie $1$.
Na $F_2$ w stanie $0$ prawdziwe jest $p$,
ale dostępny stan $1$ nie spełnia $\Diamond p$.
Stąd $\Box\Diamond p$ jest fałszywe w $0$.
Na $F_3$ w $0$ zachodzi $\Box p$,
ale stan $1$ widzi stan $2$, gdzie $p$
jest fałszywe, więc $\Box\Box p$ jest
fałszywe w $0$.
Na $F_4$ w $0$ zachodzi $\Box p$, lecz
stan $1$ widzi $0$, gdzie $p$ jest fałszywe.
Na $F_5$ zarówno $p$, jak i $\Diamond p$
są prawdziwe w $0$, lecz stan $1$ nie ma
żadnych następców i nie spełnia $\Diamond p$.
Wobec tego B oraz 5 zawodzą w $0$.
Na $F_6$ stan $2$ widzi tylko $0$, gdzie
$p$ jest prawdziwe, jednak $0$ widzi również
$1$, gdzie $p$ jest fałszywe. Zatem 4 zawodzi
w stanie $2$.

Własności relacji podane w kolumnie ważności
sprawdza się bezpośrednio: $F_0$ ma pustą
relację, $F_1$ ma wspólny jedyny następnik
$0$, $F_2$ jest łańcuchem z pętlami, a $F_3$
jest symetryczną drogą z pętlami.
Rama $F_4$ jest symetryczna i seryjna,
$F_5$ jest przechodnia, gdyż nie istnieją
dwa następujące po sobie kroki.
Na $F_6$ każdy następnik dowolnego stanu
widzi wszystkich następców tego stanu:
dla $0,1$ jest to zbiór $\{0,1\}$, a dla $2$
jest to zbiór $\{0\}$. Stąd ważność 5.
Stosujemy Lemat~\ref{lem:ncl-modal-frame-sufficient}.

Ramy $F_0,F_1$ pokazują ścisłość pierwszych
dwóch inkluzji. Rama $F_2$ oddziela T i S4
od B oraz S5, a $F_3$ oddziela T i B od S4
oraz S5. W szczególności B i S4 są
nieporównywalne.
Rama $F_4$ oddziela K od K4, a $F_5$
oddziela K od KB i K5 oraz pokazuje, że
5 nie jest twierdzeniem K4. Wreszcie
$F_6$ pokazuje, że 4 nie jest twierdzeniem K5.
\end{proof}

\begin{center}
\begin{tikzpicture}[>=Stealth,node distance=12mm]
\node (k) {$\mathrm K$};
\node[above=of k] (kd) {$\mathrm{KD}$};
\node[above=of kd] (t) {$\mathrm T$};
\node[above left=13mm and 12mm of t] (b) {$\mathrm B$};
\node[above right=13mm and 12mm of t] (s4) {$\mathrm{S4}$};
\node[above=34mm of t] (s5) {$\mathrm{S5}$};
\draw[->] (k)--(kd);
\draw[->] (kd)--(t);
\draw[->] (t)--(b);
\draw[->] (t)--(s4);
\draw[->] (b)--(s5);
\draw[->] (s4)--(s5);
\end{tikzpicture}
\end{center}
Strzałki przedstawiają ścisłą inkluzję wśród
sześciu wymienionych systemów, nie wszystkie
możliwe rozszerzenia normalnej logiki K.

\begin{stwierdzenie}[Niezależność rozdzielności operatorów]
\label{prop:ncl-modal-nondistributivity}
W K nie są dowodliwe schematy
\[
 \Box(A\vee B)\to(\Box A\vee\Box B),
 \qquad
 (\Diamond A\wedge\Diamond B)\to\Diamond(A\wedge B).
\]
Nie są one dowodliwe nawet w S5.
\end{stwierdzenie}
\begin{proof}
Niech $W=\{0,1\}$, $R=W\times W$,
$V(p)=\{0\}$ i $V(q)=\{1\}$.
Relacja jest zwrotna i spełnia warunek
schematu 5, więc wszystkie twierdzenia S5
są ważne na tej ramie.
W każdym stanie $p\vee q$ jest konieczne,
ale ani $p$, ani $q$ nie jest konieczne.
Z drugiej strony oba zdania są możliwe,
lecz $p\wedge q$ nie jest możliwe.
Obie implikacje zawodzą w każdym stanie.
\end{proof}

\begin{wniosek}[Niesprzeczność podstawowych systemów]
\label{cor:ncl-modal-basic-consistency}
Wszystkie systemy zawarte w S5 są niesprzeczne
w tym sensie, że nie dowodzą $\bot$.
\end{wniosek}
\begin{proof}
Jednoelementowa rama z pętlą spełnia warunki
T oraz 5. Wszystkie twierdzenia S5 są na niej
ważne na mocy Lematu~\ref{lem:ncl-modal-frame-sufficient},
a $\bot$ jest w jej jedynym stanie fałszywe.
\end{proof}
